\section{Study Design}\label{sec:method}

Table~\ref{tab:dataset} summarises how the dataset was built. All code is in the replication package (see the data availability statement); each step writes its counts to a file, and one script turns those files into every number, table and figure in this paper.

\subsection{Retirement registry}

We copied every retirement listed on the deprecation pages of OpenAI, Anthropic and Google into a registry: the model ID and its aliases, the announcement day, the shutdown day, the recommended replacement, and earlier plans when a retirement was postponed. For OpenAI and Anthropic the dates are exact. Google gives no announcement day and, for many models, only the earliest possible shutdown day, so we report Google separately throughout. Gemini~1.5 models, which are no longer on Google's deprecation page, are used only to tell whether a migration target was already shut down (RQ3).

\subsection{Mining candidate commits}

For each of the \ModelIds{} model IDs and aliases we ran GitHub's commit search for the exact ID between \MiningStart{} and \MiningEnd. GitHub returns at most 1,000 results per search, so each search window was split until it held fewer, down to a single day. In total we ran \Searches{} searches and saw \ResultsSeen{} results. On \TruncatedWindows{} days a single ID had more than 1,000 results and only the first 1,000 could be seen; Section~\ref{sec:robustness} shows that this does not change the results. Five IDs had more than 50,000 results in total; we searched them only from 30 days before their earliest deprecation notice, the same cut-off Kim applies. None of these five belongs to a model counted in the RQ1 headline (Section~\ref{sec:robustness}). Because GitHub's search is fuzzy, we kept only messages that contain the exact ID, dropped forks, and kept each commit once with every repository it appears in. This gave \Candidates{} candidate commits.

\subsection{Labelling migrations}\label{sec:labels}

Each candidate gets one field per rule, so any definition of a migration can be rebuilt from the same file. The rules detect a migration word (migrate, replace, switch, upgrade, remove, deprecate and similar), the direction of the message relative to the retired ID (\emph{away}, as in ``from X'', ``X to Y'', ``replace X with'', ``X is deprecated''; \emph{into}, as in ``to X'' or ``use X''; or unclear), whether the commit falls more than 30 days before the deprecation notice, dependency-bot wording, and AI coding-tool trailers. From these fields we build three labels:

\begin{itemize}
\item \texttt{migration} follows Kim's rules: a migration word, a direction other than into, not more than 30 days before the notice, not a bot, and no AI-tool trailer (\KimMigration{} commits).
\item \texttt{migration\_strict} adds that the direction must be away.
\item \texttt{migration\_away}, our main label, keeps every commit whose message moves away from a retired model and is not from a bot, whatever its wording, date or trailer (\MigrationsAway{} commits). We designed it after a first validation round, on a separate sample of 300 commits, showed that Kim's rules keep many commits that do not move away from the retired model.
\end{itemize}

\gentable{dataset}{tab:dataset}{How the dataset was built. Each line keeps the commits of the line above that pass one more rule; the last two lines are the labels.}

\subsection{Validating the label}

Because \texttt{migration\_away} was designed after looking at labelled data, we measured it on a fresh sample: 200 candidates not used before, drawn at random with a fixed seed from three strata: commits the label keeps (100), commits kept only by Kim's rules (50), and commits neither keeps (50). The labelling sheet showed only the commit and its message, never the label. The author labelled every row. An AI assistant then flagged rows that might be mislabelled, mostly commits that move into the retired model rather than away from it; the author re-checked each flagged row and made every final decision (see the declarations), and we report these final labels. A colleague independently labelled the first \FilterKappaRows{} rows from a blank copy, with no help, and agreed with $\kappa = \FilterKappa$ \citep{cohen1960}. Of the \PrecisionRows{} labelled commits the label keeps, \Precision\% are real migrations (95\% Wilson interval \PrecisionLow--\PrecisionHigh; \citealp{wilson1927}). Weighting each stratum by its size, the label finds about \Recall\% of all real migrations among the candidates. Table~\ref{tab:validation} lists every check.

\gentable{validation}{tab:validation}{Checks of the labels against hand labels. Agreement with 95\% Wilson intervals; $\kappa$ is Cohen's kappa. ``Before the fix'' is the score of the rules as they were when the sample was drawn, which is the validation result; ``after'' shows what the fix changed.}

\subsection{Placing migrations on the timeline (RQ1, RQ2)}

For each migration we take the retired model the message moves away from (if there are several, the one shut down first) and compare the commit day with that model's announcement and shutdown. A commit is \emph{after the shutdown} if it is made on or after the shutdown day. The commit day is the committer date, since GitHub search works on it.

A model shut down shortly before our data ends has had little time to collect late migrations, which would lower the share after the shutdown. Our headline therefore counts only OpenAI and Anthropic models shut down by \RetiredBy, which leaves at least 90 days of data after every shutdown, and reports the other choices next to it. For RQ2 we report the median number of days from the shutdown with a distribution-free 95\% confidence interval from order statistics \citep{conover1999}, the share of migrations in time windows around the shutdown, and, for migrations made during the notice period, how much of that period had passed.

\subsection{Finding the target model (RQ3)}

We extract the model a commit moves to from its message with eight ordered patterns, such as ``X~$\rightarrow$~Y'', ``from X to Y'', ``replace X with Y'' and ``switch to Y''. We then compare the target with the provider's recommended replacement: the exact model, a newer model of the same line (for example a newer Sonnet for a retired Sonnet), or another model. Since provider pages list what they recommend today, not on the day of the commit, the same-line measure is the fairer one and the exact match is a lower bound. We also check whether the target is itself on a retirement list, already announced for retirement, or already shut down on the day of the commit. On the hand-labelled sample the extracted target matched the hand label in \TargetCheckBefore\% of commits before we tuned the patterns on a separate, earlier sample, and in \TargetCheck\% after; when the code names a target, it names the right model in \TargetNamedRight\% of cases.

\subsection{Changed files (RQ4)}

The search results do not include changed files, so we drew a random sample of \FilesSample{} migrations (seed \FilesSeed) and fetched each commit from GitHub's API; \FilesFound{} were still available. Each changed file is sorted by its path into a kind: test, evaluation, prompt, CI, documentation, configuration, code or other. A commit \emph{touches a test} if it adds or changes a test file, and \emph{touches an evaluation} if it adds or changes an evaluation file, such as a benchmark, golden outputs or an evaluation config. A \emph{pure ID swap} changes nothing but the model ID in the lines it edits. The author checked the file kinds by hand on \FileCheckRows{} of these commits, blind to the code's answer. The rules then agreed on tests in \FileCheckTestBefore\% of commits ($\kappa = \FileCheckTestBeforeKappa$) and on evaluations in \FileCheckEvalBefore\% ($\kappa = \FileCheckEvalBeforeKappa$); the code found only \FileCheckEvalBeforeFound{} of the \FileCheckEvalHandYes{} evaluation changes seen by hand. We improved the rules from the disagreements, which raised agreement to \FileCheckTestAfter\% and \FileCheckEvalAfter\%, and we treat the evaluation share as a lower bound.
